\section{Related Work}

\paragraph{CXL synchronization.}
Suetterlein et al.~\cite{suetterleinSynchronizationCXLBased2024} explore software locks in CXL memory using a Tournament Peterson lock, comparing against single RMW atomic operations.
Their work highlights that contention spreading across multiple cache lines improves performance, a result we reproduce and extend.
However, their evaluation does not measure hardware coherence protocol overhead (they compare against bare atomic operations, not hardware locks), does not explore uncached or non-temporal access modes, and does not identify crossover points between software and hardware approaches.
Our work extends this by evaluating 13 lock implementations across four access modes with quantitative crossover analysis.
Other work has discussed the need for coherence-free synchronization in CXL~\cite{hongDawnDisaggregationCoherence2025,10.1145/3678015.3680494} but without systematic evaluation on hardware.

\paragraph{Synchronization on NUMA and distributed shared memory.}
The comprehensive study by David et al.~\cite{davidEverythingYouAlways2013} characterizes synchronization primitives across NUMA topologies, demonstrating that lock performance is architecture-dependent and that no single lock dominates across all configurations.
Our work applies a similar methodology to the CXL domain, where the additional dimension of access mode selection (cached vs.\ uncached vs.\ non-temporal) creates new trade-offs not present in traditional NUMA.
Specifically, NUMA systems always use cached coherent access; the option to bypass the cache entirely (UC/NT) is unique to CXL's disaggregated memory model.

NUMA-aware locks~\cite{diceCompactNUMAAwareLocks2019,luchangcoHierarchicalCLHQueue2006,chabbiHighPerformanceLocks2015} exploit locality through cohorting~\cite{diceLockCohortingGeneral}, where threads on the same NUMA node form cohorts to reduce cross-node traffic.
This technique is highly effective when locality is well-defined (i.e., fixed NUMA node assignments).
Adapting cohorting to CXL requires new locality abstractions: in a switched CXL fabric, ``locality'' depends on the number of switch hops between a host and the memory device, which may change dynamically as memory is remapped.

Boyd-Wickizer et al.~\cite{boyd-wickizerNonscalableLocksAre} demonstrate that non-scalable locks can cause sudden performance collapse at critical thread counts, motivating our evaluation across a wide range of concurrency levels.
Our results confirm this pattern for spin and ticket locks in CXL environments, with the collapse occurring at lower thread counts under hardware coherence due to amplified coherence traffic.

\paragraph{Lock design and fairness.}
Buhr et al.~\cite{buhrHighcontentionMutualExclusion2018} propose the Elevator lock family that we evaluate, demonstrating that software locks can achieve performance competitive with hardware locks by reducing contention through explicit successor notification.
Our CXL evaluation confirms this finding and shows that the advantage grows with CXL latency.

Beyond acquisition fairness, recent work examines \emph{usage} fairness---ensuring equitable critical section time across threads~\cite{kashyapSchedulerCooperativeLocks2020}.
In CXL environments, usage fairness interacts with access mode: under UC access, all threads pay equal CXL round-trip costs, while under Cached-HC, the thread that last held the lock benefits from cached state (the ``local spinning'' advantage).
We leave usage fairness evaluation to future work but note that UC access inherently provides more equitable costs across threads.

\paragraph{Uncached and non-temporal memory access.}
UC and NT access modes are well-studied for MMIO, persistent memory~\cite{yangOptaneMemory}, and GPU memory~\cite{wangIntelligentMultiPortMemory2008}, but have not been evaluated for inter-host synchronization over CXL.
Persistent memory work is particularly relevant: Intel Optane DIMM research extensively characterized the performance of NT stores via write-combining for persistent data structures.
Our work applies similar techniques to synchronization primitives, showing that the write-combining benefit translates to lock unlock paths.
The key difference from persistent memory is that CXL memory has higher base latency (400\,ns+ vs.\ $\sim$300\,ns for Optane) and does not require persistence guarantees, simplifying the fence requirements.
